\documentclass[../../syllabus_1.tex]{subfiles}
\begin{document}
\raggedbottom
\chapter{Polygone, cercle, périmètre et aire}

\section*{Matières}
\begin{enumerate}[label=\textbf{\arabic*.}]
    \item Le cercle et le disque.
    \item Les polygones.
    \item Formules de périmètre.    
    \item Formules d'aire.
\end{enumerate}



\newpage
\section*{Mots-clés}
\begin{enumerate}[cols=2, label=\textbf{\arabic*.}]
    \item Cercle
    \item Disque
    \item Hauteur
    \item Base
    \item Cercle inscrit
    \item Cercle circonscrit
    \item Polygone
    \item Aire
    \item Périmètre
\end{enumerate}

\section*{Attendus}
\begin{enumerate}[label=\textbf{\arabic*.}]
    \item Identifier des polygones réguliers sur base des caractéristiques : les côtés (nombre, longueur) ; les angles (amplitude).
    \item Construire une figure complexe (figure composée de figures simples), les étapes de construction étant données.
    \item Rédiger des étapes de construction d'une figure complexe (figure composée de figures simples) donnée.
    \item Construire le cercle inscrit et le cercle circonscrit à un triangle.
    \item Exprimer le périmètre d'un triangle équilatéral, d'un rectangle, d'un carré, d'un cercle en fonction de ses dimensions.
    \item Justifier les liens entre les formules d'aire du rectangle, du parallélogramme, du trapèze, du triangle et du losange.
    \item Énoncer la formule du calcul de l'aire du disque.
    \item Exprimer l'aire d'une figure simple (rectangle, parallélogramme, trapèze, triangle, losange, carré) après avoir repéré les dimensions utiles.
    \item Calculer la longueur d'un côté d'un polygone régulier (triangle équilatéral, carré, pentagone, hexagone, octogone, décagone) dont le périmètre est donné.
    \item Calculer le périmètre d'un cercle.
    \item Calculer l'aire d'une figure simple (triangle, quadrilatère, disque).
    \item Décomposer une surface complexe en plusieurs surfaces connues (triangles, quadrilatères, disques) pour en calculer l'aire, avec des formules préalablement construites.
    \item Résoudre des problèmes faisant intervenir des calculs d'aire et de périmètre de figures simples ou complexes, en situations contextualisées.
    \item Représenter une figure simple dont l'aire est identique à celle d'une autre figure simple donnée. Exemple : Représenter un rectangle dont l'aire est identique à celle d'un triangle donné.
    \item Identifier les droites remarquables d'un triangle.
    \item Définir les droites remarquables d'un triangle.
    \item Énoncer les propriétés des droites remarquables dans le triangle.
    \item Construire les droites remarquables d'un triangle ou d'un quadrilatère.
    \item Énoncer les propriétés relatives aux angles, aux axes et centre de symétrie, aux droites remarquables, dans des figures simples.
    \item Terminer la construction d'une figure simple respectant des contraintes sur les droites remarquables.
    \item Construire une figure simple dont le centre de symétrie, le(les) axe(s) de symétrie sont donnés.
\end{enumerate}

\end{document}